\section{Algorithm}
\label{sec:algorithm}

\subsection{Setup and Notation}

Let $G = (V, E)$ be the census-tract adjacency graph for a state with
$n = |V|$ tracts and $|E|$ shared boundaries.
Each vertex $v \in V$ carries a population $p(v)$.
A \emph{$k$-partition} of $G$ is a function $\sigma: V \to \{1,\ldots,k\}$
such that each part $V_i = \sigma^{-1}(i)$ induces a connected subgraph
of $G$.
A partition is \emph{population-balanced} if
\[
  \left|\frac{\sum_{v \in V_i} p(v)}{\bar{p}} - 1\right| \leq \varepsilon
  \quad \forall i \in \{1,\ldots,k\},
\]
where $\bar{p} = \frac{1}{k}\sum_{v \in V} p(v)$ is the target district
population and $\varepsilon = 0.005$ is the statutory tolerance.

\subsection{Wilson's Loop-Erased Random Walk}

\begin{definition}[Uniform Random Spanning Tree]
A \emph{uniform random spanning tree} (UST) of a connected graph $H$
is a spanning tree of $H$ drawn uniformly at random from the set of
all spanning trees of $H$.
\end{definition}

Wilson's algorithm~\citep{wilson1996generating} generates a UST of $H$
as follows:

\begin{algorithm}
\caption{Wilson's UST Algorithm}
\label{alg:wilson}
\begin{algorithmic}[1]
\Require Connected graph $H = (V_H, E_H)$, arbitrary root $r \in V_H$
\Ensure Uniform random spanning tree $T$ of $H$
\State $T \leftarrow \{r\}$
\While{$T \neq V_H$}
  \State Pick any $u \in V_H \setminus T$
  \State Perform a simple random walk from $u$ until hitting $T$, recording the path $\pi$
  \State Erase all loops from $\pi$ to obtain the loop-erased path $\mathsf{LE}(\pi)$
  \State Add all vertices and edges of $\mathsf{LE}(\pi)$ to $T$
\EndWhile
\State \Return $T$
\end{algorithmic}
\end{algorithm}

\begin{theorem}[\citealt{wilson1996generating}; planar bound \citealt{aldous1991random}]
Algorithm~\ref{alg:wilson} outputs a uniform random spanning tree of $H$
in expected time $O(\tau_{\mathrm{cover}}(H))$, where $\tau_{\mathrm{cover}}(H)$
is the cover time of a simple random walk on $H$ \citep{wilson1996generating}.
For planar graphs with $m$ vertices, Aldous~\citeyearpar{aldous1991random}
proved $\tau_{\mathrm{cover}}(H) = O(m \log m)$; this planar specialization
is not part of Wilson's original result.
\end{theorem}

\paragraph{Planarity of census-tract adjacency subgraphs.}
In practice, the subgraph $H = G[V_i \cup V_j]$ is a planar graph: it is induced from the census-tract adjacency graph, which is derived from TIGER shapefiles defining non-crossing polygonal boundaries.
Small non-planar artefacts (tri-point boundaries, state-line adjacencies) affect at most a small fraction of steps and do not alter the asymptotic cover-time bound.

In our setting, $H$ is the subgraph induced by the union $V_i \cup V_j$
of two adjacent districts chosen for recombination.
This subgraph has approximately $m = 2n/k$ vertices, giving an expected
Wilson cost of $O((n/k)\log(n/k))$ per \recom{} step.

\subsection{The ReCom Proposal}

Algorithm~\ref{alg:recom} describes one step of the \recom{} chain.
The critical design decision is the treatment of tree bipartitions:
we enumerate \emph{all} balanced cuts of the spanning tree rather than
sampling a single cut.
This enumeration is necessary to match the intended ReCom proposal family:
the implementation samples a spanning tree, enumerates every balanced tree
cut, and then chooses uniformly from the feasible cut set. The paper does not
claim that a finite run proves global uniformity over all valid plans.

In the atlas framing, a ReCom step has four inspectable candidate objects:
the adjacent district pair, the merged-region spanning tree, the balanced cut
set, and the accepted or rejected replacement pair. The chain transcript should
distinguish an accepted balanced cut from a no-balanced-cut tree, a pair
reselection event, or a slow-moving chain diagnostic. Those are different
sampling facts and support different comparison claims.

\begin{algorithm}
\caption{One ReCom Step}
\label{alg:recom}
\begin{algorithmic}[1]
\Require Current partition $\sigma$, population tolerance $\varepsilon$
\Ensure Updated partition $\sigma'$
\State $P \leftarrow$ adjacent district pairs in $\sigma$, shuffled uniformly
\For{$(i,j)$ in first 50 entries of $P$}
  \State $R \leftarrow V_i \cup V_j$ \Comment{Merged region}
  \State $H \leftarrow G[R]$ \Comment{Induced subgraph}
  \For{$r=1,\ldots,10$}
    \State $T \leftarrow \textsc{WilsonUST}(H)$ \Comment{Algorithm \ref{alg:wilson}}
    \State $\mathcal{B} \leftarrow$ all balanced tree cuts of $T$
    \If{$\mathcal{B} \neq \emptyset$}
      \State Sample $e^* \sim \mathrm{Uniform}(\mathcal{B})$
      \State $\{A, B\} \leftarrow$ components induced by removing $e^*$
      \State Update $\sigma'$ by assigning $V_i \leftarrow A$, $V_j \leftarrow B$
      \State \Return $\sigma'$
    \EndIf
  \EndFor
\EndFor
\State \Return $\sigma$ \Comment{Rejected step after pair-attempt budget}
\end{algorithmic}
\end{algorithm}

\paragraph{Failure handling.}
When a sampled spanning tree $T$ of the merged region $R$ admits no
balanced bipartition (i.e., $\mathcal{B} = \emptyset$), the algorithm
resamples the entire spanning tree from scratch rather than attempting to
repair the same tree.
This tree-resample-on-failure policy is essential: retrying with the
same tree would not sample a new point from the UST distribution and
would break the stationary distribution guarantee.
If 10 consecutive spanning tree samples for the same district pair all
fail to yield a balanced bipartition, the pair is abandoned and a new
adjacent pair $(i,j)$ is selected.
This pair-reselection mechanism is not a Rust-specific workaround; the
same failure mode occurs in Python \gc{} (Section~\ref{sec:performance}
discusses the Texas case in detail).
Rust's advantage is simply speed: it exhausts the resample budget faster
and moves on.

\paragraph{Stationarity under pair reselection.}
Pair reselection modifies the \recom{} Markov kernel: pairs with low
balanced-bipartition feasibility are abandoned more often than high-feasibility
pairs, which is not part of the original \recom{} specification.
Whether this modification preserves the stationary distribution is a
conjecture.
A plausible argument is that pair reselection is equivalent to a proposal
rejection step that discards infeasible pair proposals uniformly---which
preserves detailed balance if the pair-selection probability is symmetric
over adjacent district pairs and proposals are rejected only on balance
failure (never on cut quality).
However, this argument has not been formally verified against the specific
pair-selection distribution in Algorithm~\ref{alg:recom}.
\emph{Stationarity preservation under pair reselection is deferred to
Phase~2 empirical verification}: comparing plan-space distributions produced
by \re{} and \gc{} on the six G-track states will determine whether
any distributional difference exceeds statistical resolution.

\paragraph{Balance cut enumeration.}
For a spanning tree $T$ with $m-1$ edges, each edge removal creates a
unique bipartition of the $m$ vertices into two subtrees.
Each subtree's population is the sum of $p(v)$ over its vertices.
We scan all $m-1$ edges in a single linear pass, accumulating the
subtree population on one side via DFS, and classify each edge as
balanced or unbalanced.
The cost of this enumeration is $O(m)$, subdominant to the $O(m\log m)$
Wilson cost.

\subsection{Parallel Chains}

Independent chains are embarrassingly parallel: each chain maintains its
own partition state, its own random-number generator seeded
deterministically from the chain index, and its own step history.
No synchronization is required during sampling; chains communicate only
at the diagnostics aggregation step.

Per-chain seeds are derived deterministically.
The hash input is the fixed-width 32-byte string:
\[
  \underbrace{\texttt{b"ENSEMBLE\_CHAIN\_"}}_{15\text{ bytes}}
  \;\|\;
  \underbrace{\texttt{i.to\_le\_bytes()}}_{8\text{ bytes}}
  \;\|\;
  \underbrace{\texttt{b"\_"}}_{1\text{ byte}}
  \;\|\;
  \underbrace{\texttt{base\_seed.to\_le\_bytes()}}_{8\text{ bytes}},
\]
where $i$ is the zero-indexed chain index and \texttt{base\_seed} is
the user-supplied seed, both encoded as 8-byte little-endian unsigned
integers (\texttt{u64}).
The string literals are ASCII bytes.
The total is $15 + 8 + 1 + 8 = 32$ bytes, giving a
fixed-length, unambiguous encoding with no collision risk between chain
indices or base-seed values.
The first 8 bytes of the SHA-256 digest are interpreted as a little-endian
\texttt{u64}. The chain runner then calls
\texttt{SmallRng::seed\_from\_u64(seed)} from the \texttt{rand} crate
(version 0.8). This matches the current crate contract exactly and avoids
claiming a particular \texttt{SmallRng} internal state layout.
